% =========================================================
\section{Axisymmetric Method of Characteristics for
Supersonic Wave Analysis}
\label{sec:moc}
% =========================================================

\subsection{Motivation}

The baseline quasi-one-dimensional model provides a computationally
efficient description of the mean nozzle expansion. However, repeated
optimization runs revealed limitations that are important when the
bell fraction $K_b$ and the initial divergent wall angle
$\theta_{in}$ are treated as design variables.

The main limitations of the quasi-one-dimensional formulation are:

\begin{enumerate}
    \item it cannot resolve the propagation and interaction of
    two-dimensional expansion and compression waves;

    \item it assumes that the flow properties are uniform over each
    cross-section;

    \item it does not resolve the radial distribution of Mach number,
    pressure or flow angle at the nozzle exit;

    \item the axial-momentum loss associated with flow angularity must
    be approximated using an algebraic divergence-efficiency closure;

    \item rapid wall turning and characteristic-wave coalescence are
    not represented explicitly.
\end{enumerate}

These limitations produced the optimization behaviour discussed in
Section~\ref{sec:results}. In particular, the reduced-order model could
prefer very short contours because the viscous wall-friction penalty
was represented, while the possible inviscid benefit of smoother wave
turning and a more uniform exit flow was not resolved.

To establish a direct relationship between $\theta_{in}$, $K_b$ and
the internal supersonic flowfield, an axisymmetric Method of
Characteristics (MOC) solver is introduced. The MOC is used here as a
\emph{prescribed-wall analysis method}: the nozzle contour remains
defined by the parametric geometry model, and the MOC calculates the
flowfield inside that contour.

For each candidate geometry, the required solution variables are

\begin{equation}
    p(x,r), \qquad
    \rho(x,r), \qquad
    V(x,r), \qquad
    M(x,r), \qquad
    T(x,r), \qquad
    \theta(x,r).
\end{equation}

The MOC therefore does not replace the parametric bell-contour model.
Instead, it evaluates the inviscid performance of the contour generated
from the genetic-algorithm variables.

% ---------------------------------------------------------
\subsection{Physical interpretation of the characteristics}
% ---------------------------------------------------------

The Method of Characteristics is derived from the inviscid Euler
equations. In a steady supersonic flow, the governing equations are
hyperbolic. Consequently, information does not propagate arbitrarily
through the domain. Infinitesimal pressure disturbances propagate along
specific directions known as the \emph{characteristics} of the flow.

For a local Mach number $M>1$, the Mach angle is

\begin{equation}
    \boxed{
    \mu =
    \arcsin\left(\frac{1}{M}\right)
    }.
    \label{eq:mach_angle}
\end{equation}

If the local velocity vector forms an angle $\theta$ with the nozzle
axis, the two characteristic directions in the meridional $(x,r)$
plane are

\begin{equation}
    \boxed{
    \frac{dr}{dx}
    =
    \tan(\theta+\mu)
    }
    \qquad \text{along } C^+,
    \label{eq:cplus_direction}
\end{equation}

and

\begin{equation}
    \boxed{
    \frac{dr}{dx}
    =
    \tan(\theta-\mu)
    }
    \qquad \text{along } C^-.
    \label{eq:cminus_direction}
\end{equation}

The streamline is tangent to the local velocity direction $\theta$,
whereas the characteristics are inclined by the Mach angle $\mu$
relative to that streamline. Characteristics should therefore not be
confused with streamlines.

The intersection of one $C^+$ characteristic and one $C^-$
characteristic defines a new solution point. The position of the point
is obtained from the intersection of the characteristic directions,
while the local flow properties are obtained from the corresponding
compatibility equations.


% Replacement for the crowded four-panel Mach-angle/MOC figure.
% Required packages/libraries (already loaded by publication_manuscript/main.tex):
%   pdflscape, caption, tikz
%   \usetikzlibrary{calc,arrows.meta,angles,quotes,positioning}

\clearpage
\begin{landscape}
\thispagestyle{plain}
\vspace{-1cm}
\begin{center}
\begin{minipage}{0.485\linewidth}
\centering
\resizebox{\linewidth}{!}{%
\begin{tikzpicture}[
  line cap=round,
  line join=round,
  >=Latex,
  mainline/.style={black,line width=0.9pt},
  auxiliary/.style={red!75!black,densely dashed,line width=0.75pt},
  characteristic/.style={red!75!black,line width=1.05pt},
  knownpoint/.style={circle,fill=black,inner sep=2.2pt}
]
\path[use as bounding box] (-0.2,-0.2) rectangle (11.8,7.0);

\node[font=\bfseries\large,align=center] at (5.8,6.65)
  {(a) Geometrical origin of the Mach angle};

\coordinate (O) at (1.35,3.55);
\coordinate (D) at (6.15,3.55);
\coordinate (T) at (1.89,5.10);

\draw[red!75!black,line width=0.9pt] (O) circle (1.65);
\node[red!75!black,align=center] at (1.30,5.65)
  {acoustic wavefront\\after $\Delta t$};

\draw[mainline,->] (O) -- (D)
  node[midway,below=3pt] {$V\,\Delta t$};
\draw[auxiliary] (O) -- (T)
  node[midway,left=2pt] {$a\,\Delta t$};
\draw[characteristic] (D) -- (T);
\node[red!75!black,rotate=-20] at (4.05,4.62)
  {Mach-wave envelope};

\node[knownpoint] at (O) {};
\node[knownpoint] at (D) {};
\node[knownpoint] at (T) {};
\node[below left=1pt] at (O) {$O$};
\node[below=2pt,align=center] at (D) {disturbance\\after $\Delta t$};
\node[above=2pt] at (T) {$T$};

\draw[mainline,->] (D) -- ++(1.25,0) node[right] {$\vec V$};
\pic[draw=black,angle radius=4mm] {right angle=D--T--O};
\pic[draw=red!75!black,"$\mu$",angle eccentricity=1.45,
  angle radius=8mm] {angle=T--D--O};

\node[draw=black,rounded corners,align=center,inner sep=6pt,
  text width=3.55cm] at (9.95,3.55)
{
  The triangle $OTD$ is right-angled at $T$:\\[2mm]
  $\displaystyle
    \sin\mu=\frac{OT}{OD}
    =\frac{a\Delta t}{V\Delta t}
    =\frac{1}{M}$
  \\[2mm]
  $\displaystyle
    \boxed{\mu=\arcsin\!\left(\frac{1}{M}\right)}$
};
\end{tikzpicture}%
}
\end{minipage}
\hfill
\begin{minipage}{0.485\linewidth}
\centering
\resizebox{\linewidth}{!}{%
\begin{tikzpicture}[
  line cap=round,
  line join=round,
  >=Latex,
  mainline/.style={black,line width=0.9pt},
  characteristic/.style={red!75!black,line width=1.05pt},
  characteristicplus/.style={red!75!black,dashed,line width=1.05pt},
  knownpoint/.style={circle,fill=black,inner sep=2.2pt}
]
\path[use as bounding box] (-0.2,-0.2) rectangle (11.8,7.0);

\node[font=\bfseries\large,align=center] at (5.8,6.65)
  {(b) Characteristic directions at a flow point};
\node[align=center,text width=10.8cm] at (8.8,5.85)
  {A pressure disturbance can only propagate along the two local Mach directions.};

\coordinate (A) at (2.05,3.25);
\draw[mainline,->] (0.45,3.25) -- (6.95,3.25) node[right] {$x$};
\node[knownpoint] at (A) {};
\node[below left=2pt] at (A) {$A$};

\draw[mainline,->] (A) -- ++(18:3.45) node[right] {$\vec V_A$};
\draw[characteristicplus,->] (A) -- ++(48:3.30)
  node[above right] {$C^+$};
\draw[characteristic,->] (A) -- ++(-12:3.65)
  node[below right] {$C^-$};

\draw[black] ($(A)+(0.72,0)$)
  arc[start angle=0,end angle=18,radius=0.72];
\node at ($(A)+(0.96,0.18)$) {$\theta_A$};

\draw[red!75!black] ($(A)+(18:1.08)$)
  arc[start angle=18,end angle=48,radius=1.08];
\node[red!75!black] at ($(A)+(35:1.40)$) {$\mu_A$};

\draw[red!75!black] ($(A)+(18:0.96)$)
  arc[start angle=18,end angle=-12,radius=0.96];
\node[red!75!black] at ($(A)+(1:1.30)$) {$\mu_A$};

\node[draw=black,rounded corners,align=left,inner sep=7pt,
  text width=3.70cm] at (9.65,3.15)
{
  Flow direction:\quad $\theta_A$\\[2mm]
  Upper characteristic:\quad $\theta_A+\mu_A$\\[2mm]
  Lower characteristic:\quad $\theta_A-\mu_A$
};
\end{tikzpicture}%
}
\end{minipage}

\vspace{0.55cm}

\begin{minipage}{0.485\linewidth}
\centering
\resizebox{\linewidth}{!}{%
\begin{tikzpicture}[
  line cap=round,
  line join=round,
  >=Latex,
  mainline/.style={black,line width=0.9pt},
  auxiliary/.style={red!75!black,densely dashed,line width=0.75pt},
  characteristic/.style={red!75!black,line width=1.05pt},
  knownpoint/.style={circle,fill=black,inner sep=2.2pt},
  newpoint/.style={circle,draw=black,fill=white,line width=1pt,inner sep=3.2pt}
]
\path[use as bounding box] (-0.2,-0.2) rectangle (11.8,7.0);

\node[font=\bfseries\large,align=center] at (5.8,6.65)
  {(c) Point--slope equation of one characteristic};
\node[align=center,text width=10.8cm] at (5.8,5.88)
  {The horizontal and vertical legs are coordinate differences, not new flow properties.};

\draw[mainline,->] (0.45,1.00) -- (6.65,1.00) node[right] {$x$};
\draw[mainline,->] (0.70,0.70) -- (0.70,5.45) node[above] {$r$};

\coordinate (A) at (1.45,4.80);
\coordinate (P) at (5.80,1.75);
\coordinate (H) at (5.80,4.80);

\draw[characteristic,->] (A) -- (P)
  node[midway,below,sloped] {$C^-$ characteristic};
\draw[auxiliary] (A) -- (H);
\draw[auxiliary] (H) -- (P);

\draw[gray,densely dotted] (A) -- (1.45,1.00);
\draw[gray,densely dotted] (P) -- (5.80,1.00);
\draw[gray,densely dotted] (A) -- (0.70,4.80);
\draw[gray,densely dotted] (P) -- (0.70,1.75);

\node[knownpoint] at (A) {};
\node[newpoint] at (P) {};
\node[knownpoint] at (H) {};
\node[above left=2pt] at (A) {$A=(x_A,r_A)$};
\node[below right=2pt] at (P) {$(x,r)$};
\node[above right=1pt] at (H) {$H=(x,r_A)$};

\draw[<->,black,line width=0.8pt] ($(A)+(0,-0.40)$) --
  ($(H)+(0,-0.40)$)
  node[midway,fill=white,inner sep=2pt] {$x-x_A$};
\draw[<->,black,line width=0.8pt] ($(H)+(0.42,0)$) --
  ($(P)+(0.42,0)$)
  node[midway,right=3pt,align=left] {$r-r_A$\\[-1mm]$(<0)$};

\draw[red!75!black] ($(A)+(0.78,0)$)
  arc[start angle=0,end angle=-35.0,radius=0.78];
\node[red!75!black] at ($(A)+(0.92,-0.58)$) {$\alpha_A$};

\node[draw=black,rounded corners,align=center,inner sep=7pt,
  text width=4.15cm] at (9.15,3.10)
{
  $\displaystyle
  \text{slope}
  =\frac{\text{vertical change}}{\text{horizontal change}}
  =\frac{r-r_A}{x-x_A}
  =\tan\alpha_A$
  \\[2mm]
  Since $\alpha_A=\theta_A-\mu_A$,\\[1mm]
  $\displaystyle
  \boxed{r-r_A=\tan(\theta_A-\mu_A)(x-x_A)}$
};
\end{tikzpicture}%
}
\end{minipage}
\hfill
\begin{minipage}{0.485\linewidth}
\centering
\resizebox{\linewidth}{!}{%
\begin{tikzpicture}[
  line cap=round,
  line join=round,
  >=Latex,
  mainline/.style={black,line width=0.9pt},
  characteristic/.style={red!75!black,line width=1.05pt},
  characteristicplus/.style={red!75!black,dashed,line width=1.05pt},
  knownpoint/.style={circle,fill=black,inner sep=2.2pt},
  newpoint/.style={circle,draw=black,fill=white,line width=1pt,inner sep=3.2pt}
]
\path[use as bounding box] (-0.2,-0.2) rectangle (11.8,7.0);

\node[font=\bfseries\large,align=center] at (5.8,6.65)
  {(d) Creation of a new MOC point};

\draw[gray,densely dotted,line width=0.9pt]
  (1.00,1.10) -- (1.00,5.55);
\node[rotate=90] at (0.68,3.30) {known line};

\coordinate (A) at (1.00,5.05);
\coordinate (B) at (1.00,1.60);
\coordinate (P) at (5.40,3.25);

\node[knownpoint] at (A) {};
\node[knownpoint] at (B) {};
\node[newpoint] at (P) {};
\node[left=3pt] at (A) {$A$};
\node[left=3pt] at (B) {$B$};
\node[above right=2pt] at (P) {$P=(x_P,r_P)$};

\draw[characteristic,->] (A) -- (P)
  node[midway,above,sloped] {$C^-$};
\draw[characteristicplus,->] (B) -- (P)
  node[midway,below,sloped] {$C^+$};
\draw[mainline,->] (P) -- ++(1.05,0.20) node[right] {$\vec V_P$};

\node[red!75!black,align=center] at (3.05,5.05)
  {$m_A=\tan(\theta_A-\mu_A)$};
\node[red!75!black,align=center] at (3.05,1.48)
  {$m_B=\tan(\theta_B+\mu_B)$};

\node[draw=black,rounded corners,align=center,inner sep=7pt,
  text width=4.35cm] at (9.05,3.35)
{
  The two characteristic equations are
  \[
    r-r_A=m_A(x-x_A),\qquad
    r-r_B=m_B(x-x_B).
  \]
  At the intersection, both have the same $x=x_P$ and $r=r_P$:
  \[
    \boxed{x_P=
    \frac{r_B-r_A+m_Ax_A-m_Bx_B}{m_A-m_B}}
  \]
  \vspace{-2mm}
  \[
    \boxed{r_P=r_A+m_A(x_P-x_A)}.
  \]
};
\end{tikzpicture}%
}
\end{minipage}

\vspace{0.35cm}

\captionof{figure}{%
Geometrical interpretation of the Mach angle and construction of a new
Method-of-Characteristics point. 
}
\label{fig:mach_angle_characteristic_geometry}
\end{center}

\end{landscape}
\clearpage
% ---------------------------------------------------------
\subsection{Characteristic compatibility equations}
% ---------------------------------------------------------

The characteristic directions alone determine where information
propagates, but they do not determine how the flow variables change.
The required additional relations are the
\emph{compatibility equations}.

A compatibility equation expresses the combinations of pressure,
velocity and flow angle that must be satisfied between two points
belonging to the same characteristic.

A convenient pressure-based axisymmetric formulation is

\begin{equation}
    Q\,dp+d\theta+S_+\,dx=0
    \qquad \text{along } C^+,
    \label{eq:compatibility_plus}
\end{equation}

and

\begin{equation}
    Q\,dp-d\theta+S_-\,dx=0
    \qquad \text{along } C^-,
    \label{eq:compatibility_minus}
\end{equation}

where

\begin{equation}
    Q
    =
    \frac{\sqrt{M^2-1}}{\rho V^2},
    \label{eq:moc_Q}
\end{equation}

and

\begin{equation}
    S_\pm
    =
    \frac{\sin\theta}
    {rM\cos(\theta\pm\mu)}.
    \label{eq:moc_axisymmetric_source}
\end{equation}

The terms $S_+$ and $S_-$ represent the axisymmetric correction caused
by the changing circumference of the streamtube. They vanish in the
planar-flow limit.

For a planar, isentropic and calorically perfect flow, the compatibility
equations reduce to the familiar Prandtl--Meyer invariants

\begin{equation}
    d(\theta-\nu)=0
    \qquad \text{along } C^+,
\end{equation}

and

\begin{equation}
    d(\theta+\nu)=0
    \qquad \text{along } C^-,
\end{equation}

where $\nu=\nu(M)$ is the Prandtl--Meyer function:
\begin{equation}
\nu(M)
=
\sqrt{\frac{\gamma+1}{\gamma-1}}
\tan^{-1}
\left[
\sqrt{
\frac{\gamma-1}{\gamma+1}
\left(M^2-1\right)
}
\right]
-
\tan^{-1}
\left(
\sqrt{M^2-1}
\right).
\label{eq:prandtl_meyer}
\end{equation}
In the axisymmetric formulation these quantities are no longer exactly
constant, and the geometric source terms must be integrated numerically
along each characteristic segment.

The sign convention used in
Equations~\eqref{eq:cplus_direction}--\eqref{eq:compatibility_minus}
must remain consistent throughout the implementation. Alternative
references may interchange the names $C^+$ and $C^-$.

% ---------------------------------------------------------
\subsection{Characteristic mesh and numerical marching}
% ---------------------------------------------------------

The MOC does not use a conventional rectangular computational grid.
Instead, the numerical mesh is generated by the intersections of the
$C^+$ and $C^-$ characteristic families.

The solution begins from a transonic initial line located immediately
downstream of the geometric . This line contains a discrete
distribution of the form

\begin{equation}
    \left\{
    x_i,\,
    r_i,\,
    p_i,\,
    \rho_i,\,
    V_i,\,
    M_i,\,
    \theta_i
    \right\},
    \qquad i=1,\ldots,N.
\end{equation}

Because the characteristic equations become degenerate at $M=1$, a
uniform sonic line cannot be used directly. The initial data should be
obtained from an axisymmetric transonic throat solution based on the
local throat curvature.

For an interior point $P$, one $C^+$ characteristic is traced from a
known point $A$ and one $C^-$ characteristic from a known point $B$.
Their approximate intersection provides the new coordinates
$(x_P,r_P)$. The two compatibility equations then provide the local
pressure and flow angle.

Since the characteristic slopes depend on the unknown properties at
$P$, the solution is obtained using a predictor--corrector procedure:

\begin{enumerate}
    \item calculate a first estimate using the properties at $A$ and
    $B$;

    \item evaluate the predicted properties at $P$;

    \item recompute the characteristic slopes using averaged properties;

    \item repeat until the changes in position and flow properties fall
    below the prescribed convergence tolerance.
\end{enumerate}

Interpolation is required whenever the foot of a new characteristic
falls between two nodes of the previous characteristic line. The resulting characteristic mesh and the local construction of an
interior solution point are illustrated schematically in
Figure~\ref{fig:prescribed_wall_moc_mesh}.

\begin{figure}[htbp]
\centering

% =========================================================
% (a) Global characteristic mesh
% =========================================================

\resizebox{0.98\textwidth}{!}{%
\begin{tikzpicture}[
    x=1.25cm,
    y=1.25cm,
    >=Latex,
    wall/.style={very thick,draw=black},
    axisline/.style={thick,draw=black},
    cplus/.style={thin,dashed,draw=red!75!black},
    cminus/.style={thin,draw=red!75!black},
    initial/.style={thick,double,draw=black},
    point/.style={circle,fill=black,inner sep=1.7pt}
]

% ---------------------------------------------------------
% Coordinate system
% ---------------------------------------------------------

\draw[axisline,->]
    (-1.25,0) -- (8.65,0)
    node[right] {$x$};

\draw[axisline,->]
    (0,0) -- (0,4.65)
    node[above] {$r$};

\node[below] at (4.0,-0.14)
    {axis of symmetry: $\theta=0$};

% ---------------------------------------------------------
% Prescribed nozzle wall
% ---------------------------------------------------------

% Convergent section
\draw[wall]
    (-1.15,2.65)
    --
    (-0.35,1.25);

% Throat arc and bell contour
\draw[wall]
    plot[smooth] coordinates {
        (-0.35,1.25)
        (-0.12,1.04)
        ( 0.00,1.00)
        ( 0.25,1.08)
        ( 0.65,1.30)
        ( 1.15,1.55)
        ( 2.00,1.85)
        ( 3.00,2.12)
        ( 4.20,2.40)
        ( 5.50,2.66)
        ( 6.75,2.86)
        ( 8.00,3.00)
    };

\node[above left] at (-0.65,1.56)
    {convergent};

\node[above] at (4.45,2.78)
    {prescribed parabolic contour $r_w(x)$};

% ---------------------------------------------------------
% Throat and exit
% ---------------------------------------------------------

\draw[densely dotted]
    (0,0) -- (0,1.00);

\node[point] at (0,1.00) {};
\node[above left] at (0,0.500)
    {throat};

\draw[densely dotted]
    (8.00,0) -- (8.00,3.00);

\node[above] at (8.00,3.05)
    {exit plane};

% ---------------------------------------------------------
% Initial transonic line
% ---------------------------------------------------------

\draw[initial]
    (0.24,1.075) -- (0.58,0);


% ---------------------------------------------------------
% C-minus characteristics
% Solid red lines: wall towards axis
% ---------------------------------------------------------

\draw[cminus]
    (0.24,1.075) -- (1.35,0);

\draw[cminus]
    (1.15,1.55) -- (2.45,0);

\draw[cminus]
    (2.00,1.85) -- (3.65,0);

\draw[cminus]
    (3.00,2.12) -- (4.90,0);

\draw[cminus]
    (4.20,2.40) -- (6.20,0);

\draw[cminus]
    (5.50,2.66) -- (7.55,0);

\draw[cminus]
    (6.75,2.86) -- (8.00,1.55);

% ---------------------------------------------------------
% C-plus characteristics
% Dashed red lines: axis towards wall
% ---------------------------------------------------------

\draw[cplus]
    (0.58,0) -- (1.15,1.55);

\draw[cplus]
    (1.35,0) -- (2.00,1.85);

\draw[cplus]
    (2.45,0) -- (3.00,2.12);

\draw[cplus]
    (3.65,0) -- (4.20,2.40);

\draw[cplus]
    (4.90,0) -- (5.50,2.66);

\draw[cplus]
    (6.20,0) -- (6.75,2.86);

\draw[cplus]
    (7.55,0) -- (8.00,2.45);

% ---------------------------------------------------------
% Wall tangency
% ---------------------------------------------------------

\node[point] (W) at (4.20,2.40) {};

\draw[->,thick]
    (W) -- ++(0.90,0.18)
    node[below] {$\vec V_w$};

\node[
    fill=white,
    inner sep=1pt
] at (4.2,2.07)
    {$\theta=\theta_w$};

% ---------------------------------------------------------
% Flow direction
% ---------------------------------------------------------

\draw[->,very thick]
    (6.30,1.20) -- (7.35,1.20)
    node[midway,above] {flow};

% ---------------------------------------------------------
% Legend
% ---------------------------------------------------------

\begin{scope}[shift={(0.15,3.38)}]

    \draw[cminus]
        (0,1) -- (0.40,1);

    \node[right] at (0.45,1)
        {$C^-$ characteristic};

    \draw[cplus]
        (2.90,1) -- (3.30,1);

    \node[right] at (3.35,1)
        {$C^+$ characteristic};

    \draw[initial]
        (5.70,1) -- (6.10,1);
        
    \node[right] at (6.15,1)
        {transonic initial line};
\end{scope}

% Panel identifier
\node[anchor=north west,font=\bfseries]
    at (-1.15,-0.42) {(a)};

\end{tikzpicture}%
}

\vspace{0.7cm}

% =========================================================
% (b) Close-up of one characteristic intersection
% =========================================================

\resizebox{0.76\textwidth}{!}{%
\begin{tikzpicture}[
    x=1.45cm,
    y=1.25cm,
    >=Latex,
    cplus/.style={very thick,dashed,draw=red!75!black},
    cminus/.style={very thick,draw=red!75!black},
    knownpoint/.style={
        circle,
        fill=black,
        inner sep=2.3pt
    },
    newpoint/.style={
        circle,
        draw=black,
        very thick,
        fill=white,
        inner sep=3.2pt
    }
]

% ---------------------------------------------------------
% Previous characteristic/data line
% ---------------------------------------------------------

\draw[thick,densely dotted]
    (0,0.35) -- (0,3.20);

\node[
    rotate=90,
    above,
    align=center
]
at (-0.12,1.75)
{\footnotesize previous known\\line};

% ---------------------------------------------------------
% Known points A and B
% ---------------------------------------------------------

\node[knownpoint] (A) at (0,2.75) {};
\node[left=3pt of A]
    {$A$};

\node[knownpoint] (B) at (0,0.75) {};
\node[left=3pt of B]
    {$B$};

\node[align=right,anchor=east]
    at (-0.20,2.75)
{
$\left(x_A,r_A,p_A,M_A,\theta_A\right)$
};

\node[align=right,anchor=east]
    at (-0.20,0.75)
{
$\left(x_B,r_B,p_B,M_B,\theta_B\right)$
};

% ---------------------------------------------------------
% New point P
% ---------------------------------------------------------

\node[newpoint] (P) at (3.00,1.65) {};

\node[above right=2pt and 3pt of P]
    {$P$};

% C-minus from A
\draw[cminus,->]
    (A) -- (P)
    node[midway,above,sloped]
    {$C^-$};

% C-plus from B
\draw[cplus,->]
    (B) -- (P)
    node[midway,below,sloped]
    {$C^+$};

% ---------------------------------------------------------
% Local characteristic directions
% ---------------------------------------------------------

\node[
    align=center,
    anchor=south
] at (1.35,2.55)
{
slope:
$\tan(\theta_A-\mu_A)$
};

\node[
    align=center,
    anchor=north
] at (1.30,0.78)
{
slope:
$\tan(\theta_B+\mu_B)$
};

% ---------------------------------------------------------
% State calculation at P
% ---------------------------------------------------------

\draw[->,thick]
    (3.15,1.65) -- (4.05,1.65);

\node[
    draw=black,
    rounded corners,
    align=left,
    anchor=west,
    inner sep=5pt
] at (4.08,1.65)
{
\textbf{New solution point}\\[1mm]
Geometry:
$(x_P,r_P)$\\[1mm]
Compatibility:
$(p_P,\theta_P)$\\[1mm]
Thermodynamics:
$(\rho_P,V_P,M_P,T_P)$
};

% ---------------------------------------------------------
% Compatibility labels
% ---------------------------------------------------------

\node[
    align=center,
    fill=white,
    inner sep=2pt
] at (1.500,3.87)
{
compatibility transported\\
from point $A$
};

\node[
    align=center,
    fill=white,
    inner sep=2pt
] at (1.500,-0.72)
{
compatibility transported\\
from point $B$
};

% ---------------------------------------------------------
% Flow vector at P
% ---------------------------------------------------------

\draw[->,thick]
    (P) -- ++(0.85,0.18)
    node[above] {$\vec V_P$};

% Panel identifier
\node[anchor=north west,font=\bfseries]
    at (-1.85,0.18) {(b)};

\end{tikzpicture}%
}

\caption{
Schematic representation of the axisymmetric Method of
Characteristics used for prescribed-wall nozzle analysis.
(a) Characteristic mesh inside the parametric bell-nozzle contour.
The solid and dashed red lines represent the $C^-$ and $C^+$
characteristic families, respectively.
(b) Local construction of a new interior point $P$ from the
intersection of two characteristics originating from the known
points $A$ and $B$.
}

\label{fig:prescribed_wall_moc_mesh}

\end{figure}

As shown in Figure~\ref{fig:prescribed_wall_moc_mesh}(b), the
coordinates $(x_P,r_P)$ are obtained from the intersection of the
two characteristic directions. The local pressure and flow angle are
then determined from the two compatibility equations. The remaining
thermodynamic properties follow from the energy equation and the
selected gas model.

At the nozzle wall, the flow-tangency boundary condition is

\begin{equation}
    \boxed{
    \theta_w
    =
    \arctan\left(\frac{dr_w}{dx}\right)
    }.
    \label{eq:moc_wall_condition}
\end{equation}

At the axis of symmetry,

\begin{equation}
    \boxed{
    r=0,
    \qquad
    \theta=0
    }.
    \label{eq:moc_axis_condition}
\end{equation}

The $1/r$ dependence of the axisymmetric source term requires a
dedicated limiting formulation at the axis. Direct evaluation of
Equation~\eqref{eq:moc_axisymmetric_source} at $r=0$ is not permitted.

The resulting characteristic mesh is illustrated in
Figure~\ref{fig:prescribed_wall_moc_mesh}.

% Insert the monochromatic TikZ figure here.


\subsection{Axis-of-symmetry limiting treatment}
\label{sec:moc_axis_treatment}

The axis of symmetry requires a dedicated boundary formulation. The
axisymmetric compatibility source terms,

\begin{equation}
S^\pm =
\frac{\sin\theta}
{rM\cos(\theta\pm\mu)},
\label{eq:moc_axisymmetric_source_recalled}
\end{equation}

cannot be evaluated directly at \(r=0\), because the symmetry condition
also imposes \(\theta=0\). Direct substitution would therefore produce
the indeterminate form \(0/0\). This is a coordinate singularity rather
than a physical singularity, and its finite limit follows from the
regularity and parity of the flow variables near the axis.

\subsubsection{Parity of the flow variables}

To examine the local behaviour around the axis, the radial coordinate
may temporarily be extended to a signed coordinate. Reflection across
the axis corresponds to

\begin{equation}
r\longrightarrow-r.
\end{equation}

The scalar flow properties are unchanged by this reflection:

\begin{equation}
p(x,-r)=p(x,r),
\qquad
\rho(x,-r)=\rho(x,r),
\end{equation}

\begin{equation}
T(x,-r)=T(x,r),
\qquad
M(x,-r)=M(x,r).
\end{equation}

The axial velocity component \(u\) is also unchanged,

\begin{equation}
u(x,-r)=u(x,r),
\end{equation}

whereas the radial velocity component \(v\) reverses direction:

\begin{equation}
v(x,-r)=-v(x,r).
\end{equation}

Consequently, \(u\) is locally an even function of \(r\), while \(v\)
is locally an odd function. For a smooth solution, their Taylor
expansions near \(r=0\) are therefore

\begin{equation}
u(x,r)
=
u_0(x)+u_2(x)r^2
+\mathcal{O}(r^4),
\label{eq:axis_u_expansion}
\end{equation}

and

\begin{equation}
v(x,r)
=
v_1(x)r+v_3(x)r^3
+\mathcal{O}(r^5).
\label{eq:axis_v_expansion}
\end{equation}

Since the local flow angle is defined by

\begin{equation}
\theta(x,r)
=
\tan^{-1}
\left(
\frac{v(x,r)}{u(x,r)}
\right),
\label{eq:flow_angle_velocity_definition}
\end{equation}

it is also an odd function of the radial coordinate:

\begin{equation}
\theta(x,-r)=-\theta(x,r).
\end{equation}

Its local Taylor expansion consequently contains only odd powers:

\begin{equation}
\boxed{
\theta(x,r)
=
c_1(x)r+c_3(x)r^3
+\mathcal{O}(r^5)
}.
\label{eq:axis_theta_expansion}
\end{equation}

The coefficients in Equation~\eqref{eq:axis_theta_expansion} are not
empirical loss coefficients. They are local derivatives of the
axisymmetric flowfield. In particular,

\begin{equation}
\boxed{
c_1(x)
=
\left.
\frac{\partial\theta}{\partial r}
\right|_{r=0}
}.
\label{eq:c1_axis_derivative}
\end{equation}

Expanding Equation~\eqref{eq:flow_angle_velocity_definition} using
Equations~\eqref{eq:axis_u_expansion} and
\eqref{eq:axis_v_expansion} gives

\begin{equation}
c_1(x)=\frac{v_1(x)}{u_0(x)},
\end{equation}

and

\begin{equation}
c_3(x)
=
\frac{v_3(x)}{u_0(x)}
-
\frac{v_1(x)u_2(x)}{u_0(x)^2}
-
\frac{1}{3}
\left(
\frac{v_1(x)}{u_0(x)}
\right)^3.
\end{equation}

Equation~\eqref{eq:axis_theta_expansion} is only a local expansion
around the axis. It does not assume that the complete radial flow-angle
distribution is a cubic polynomial.

\subsubsection{Finite limit of the axisymmetric source term}

Using Equation~\eqref{eq:axis_theta_expansion}, the flow angle close to
the axis satisfies

\begin{equation}
\sin\theta
=
c_1r+\mathcal{O}(r^3).
\end{equation}

It follows that

\begin{equation}
\lim_{r\rightarrow0}
\frac{\sin\theta}{r}
=
c_1
=
\left.
\frac{\partial\theta}{\partial r}
\right|_{r=0}.
\label{eq:sin_theta_over_r_limit}
\end{equation}

At the axis, \(\theta=0\), and therefore

\begin{equation}
\cos(\theta\pm\mu)
\longrightarrow
\cos(\pm\mu)
=
\cos\mu.
\end{equation}

The limiting value of Equation~
\eqref{eq:moc_axisymmetric_source_recalled} is consequently

\begin{equation}
\boxed{
S^\pm_{\mathrm{axis}}
=
\frac{1}
{M_{\mathrm{axis}}\cos\mu_{\mathrm{axis}}}
\left.
\frac{\partial\theta}{\partial r}
\right|_{r=0}
}.
\label{eq:S_axis_first_form}
\end{equation}

Since

\begin{equation}
\sin\mu_{\mathrm{axis}}
=
\frac{1}{M_{\mathrm{axis}}},
\end{equation}

then

\begin{equation}
M_{\mathrm{axis}}\cos\mu_{\mathrm{axis}}
=
\sqrt{M_{\mathrm{axis}}^2-1}.
\end{equation}

The axis limit can therefore also be written as

\begin{equation}
\boxed{
S^\pm_{\mathrm{axis}}
=
\frac{1}
{\sqrt{M_{\mathrm{axis}}^2-1}}
\left.
\frac{\partial\theta}{\partial r}
\right|_{r=0}
}.
\label{eq:S_axis_final}
\end{equation}

The two characteristic families have the same limiting source term at
the axis because

\begin{equation}
\cos(+\mu_{\mathrm{axis}})
=
\cos(-\mu_{\mathrm{axis}}).
\end{equation}

No artificial minimum radius is therefore required. Replacing \(r\)
with \(\max(r,\varepsilon_r)\) would alter the compatibility equations
and conceal the coordinate singularity instead of resolving it.

\subsubsection{Construction of a characteristic point on the axis}

Consider a known interior point

\begin{equation}
A=
\left(
x_A,r_A,p_A,M_A,\theta_A
\right),
\qquad r_A>0,
\end{equation}

from which a \(C^-\) characteristic reaches the axis at the new point

\begin{equation}
P=
\left(
x_P,0,p_P,M_P,0
\right).
\end{equation}

The axis conditions provide

\begin{equation}
\boxed{
r_P=0,
\qquad
\theta_P=0
}.
\label{eq:moc_axis_conditions}
\end{equation}

The \(C^-\) characteristic direction is

\begin{equation}
\frac{dr}{dx}
=
\tan(\theta-\mu).
\end{equation}

Using the properties at \(A\) as a first-order geometrical predictor,

\begin{equation}
-r_A
=
(x_P-x_A)
\tan(\theta_A-\mu_A),
\end{equation}

which gives

\begin{equation}
\boxed{
x_P
=
x_A
-
\frac{r_A}
{\tan(\theta_A-\mu_A)}
}.
\label{eq:axis_point_predictor}
\end{equation}

The pressure-based \(C^-\) compatibility equation is

\begin{equation}
Q\,dp-d\theta+S^-\,dx=0.
\label{eq:Cminus_axis_compatibility}
\end{equation}

A second-order finite-difference approximation between \(A\) and \(P\)
is

\begin{equation}
\overline{Q}_{AP}(p_P-p_A)
-
(\theta_P-\theta_A)
+
\overline{S^-}_{AP}(x_P-x_A)
=
0,
\label{eq:axis_discrete_compatibility}
\end{equation}

where

\begin{equation}
\overline{Q}_{AP}
=
\frac{Q_A+Q_P}{2},
\end{equation}

and

\begin{equation}
\overline{S^-}_{AP}
=
\frac{S^-_A+S^-_{\mathrm{axis}}}{2}.
\end{equation}

Using \(\theta_P=0\), the axis pressure is obtained from

\begin{equation}
\boxed{
p_P
=
p_A
-
\frac{
\theta_A+
\overline{S^-}_{AP}(x_P-x_A)
}{
\overline{Q}_{AP}
}
}.
\label{eq:axis_pressure_update}
\end{equation}

Because \(Q_P\), \(S^-_{\mathrm{axis}}\), \(\mu_P\) and the corrected
characteristic slope depend on the initially unknown state at \(P\),
Equations~\eqref{eq:axis_point_predictor}--
\eqref{eq:axis_pressure_update} are solved using the same
predictor--corrector strategy employed for an interior characteristic
point:

\begin{enumerate}
    \item predict \(x_P\) using the properties at \(A\);
    \item impose \(r_P=0\) and \(\theta_P=0\);
    \item predict \(p_P\) from the \(C^-\) compatibility equation;
    \item calculate the remaining thermodynamic properties and obtain
    \(M_P\) and \(\mu_P\);
    \item evaluate the limiting source
    \(S^-_{\mathrm{axis}}\);
    \item recompute the characteristic geometry and averaged
    compatibility coefficients;
    \item repeat until both the position and flow properties converge.
\end{enumerate}

For the corrected geometry, the characteristic slope is averaged
between \(A\) and \(P\):

\begin{equation}
-r_A
=
\frac{x_P-x_A}{2}
\left[
\tan(\theta_A-\mu_A)
+
\tan(\theta_P-\mu_P)
\right].
\end{equation}

Since \(\theta_P=0\),

\begin{equation}
\boxed{
-r_A
=
\frac{x_P-x_A}{2}
\left[
\tan(\theta_A-\mu_A)
-
\tan\mu_P
\right]
}.
\label{eq:axis_corrected_geometry}
\end{equation}

\subsubsection{Numerical estimation of the axis derivative}

The derivative required by Equation~\eqref{eq:S_axis_final} should be
estimated from interior points at the same axial characteristic station
as the new axis point. Using the closest interior point \(Q_1\),

\begin{equation}
\left.
\frac{\partial\theta}{\partial r}
\right|_{r=0}
\approx
\frac{\theta_{Q_1}}{r_{Q_1}},
\label{eq:axis_derivative_first_order}
\end{equation}

provides a first-order estimate.

A higher-order estimate can be obtained from two interior points
\(Q_1\) and \(Q_2\) by fitting the symmetry-preserving expansion

\begin{equation}
\theta(r)=c_1r+c_3r^3.
\end{equation}

The coefficients satisfy

\begin{equation}
\begin{bmatrix}
r_{Q_1} & r_{Q_1}^3\\
r_{Q_2} & r_{Q_2}^3
\end{bmatrix}
\begin{bmatrix}
c_1\\
c_3
\end{bmatrix}
=
\begin{bmatrix}
\theta_{Q_1}\\
\theta_{Q_2}
\end{bmatrix}.
\label{eq:axis_odd_polynomial_fit}
\end{equation}

The required axis derivative is then

\begin{equation}
\boxed{
\left.
\frac{\partial\theta}{\partial r}
\right|_{r=0}
=
c_1
=
\frac{
r_{Q_2}^{\,2}\theta_{Q_1}/r_{Q_1}
-
r_{Q_1}^{\,2}\theta_{Q_2}/r_{Q_2}
}{
r_{Q_2}^{\,2}-r_{Q_1}^{\,2}
}
}.
\label{eq:axis_derivative_second_order}
\end{equation}

Using a point located at a different axial coordinate, such as the
incoming characteristic point \(A\), through the approximation
\(\theta_A/r_A\), is acceptable only as an initial predictor. The
converged axis derivative should be obtained from points belonging to
the same downstream characteristic station.
% ---------------------------------------------------------
\subsection{Exit-plane thrust integration}
% ---------------------------------------------------------

The MOC provides a non-uniform radial flow profile at the nozzle exit.
The local axial velocity is

\begin{equation}
    u_e(r)=V_e(r)\cos\theta_e(r).
\end{equation}

The exit mass flow rate is evaluated from

\begin{equation}
    \boxed{
    \dot{m}_e
    =
    2\pi
    \int_0^{R_e}
    \rho_e(r)u_e(r)\,r\,dr
    }.
    \label{eq:moc_exit_massflow}
\end{equation}

The inviscid thrust is

\begin{equation}
    \boxed{
    F_{\mathrm{MOC}}
    =
    2\pi
    \int_0^{R_e}
    \left[
        \rho_e(r)u_e^2(r)
        +
        p_e(r)-p_a
    \right]r\,dr
    }.
    \label{eq:moc_thrust}
\end{equation}

The corresponding inviscid thrust coefficient is

\begin{equation}
    \boxed{
    C_{F,\mathrm{MOC}}
    =
    \frac{F_{\mathrm{MOC}}}{p_cA_t}
    }.
    \label{eq:moc_thrust_coefficient}
\end{equation}

Because the local velocity direction is resolved over the complete exit
plane, the axial-momentum effect of flow divergence is already included
in Equation~\eqref{eq:moc_thrust}. Consequently, a separate algebraic
factor such as $\eta_{div}=\cos\theta_{out}$ is not required in the
MOC-based evaluation.

% ---------------------------------------------------------
\subsection{Coupling with the boundary-layer model}
% ---------------------------------------------------------

The inviscid MOC solution supplies the boundary-layer edge conditions

\begin{equation}
    p_e(x),\qquad
    T_e(x),\qquad
    \rho_e(x),\qquad
    V_e(x),\qquad
    M_e(x)
\end{equation}

along the nozzle wall. These quantities are passed to the integral
boundary-layer model, which calculates

\begin{equation}
    \tau_w(x),\qquad
    C_f(x),\qquad
    \delta^*(x).
\end{equation}

The axial wall-friction force is

\begin{equation}
    \boxed{
    F_f
    =
    2\pi
    \int_{x_t}^{x_e}
    \tau_w(x)r_w(x)\,dx
    }.
    \label{eq:moc_wall_friction}
\end{equation}

The loss-corrected thrust coefficient is therefore

\begin{equation}
    \boxed{
    C_{F,\mathrm{eff}}
    =
    \frac{F_{\mathrm{MOC}}-F_f}{p_cA_t}
    }.
    \label{eq:moc_effective_cf}
\end{equation}

Boundary-layer separation is not predicted by the inviscid MOC itself.
Instead, the MOC provides the wall-pressure distribution and adverse
pressure gradient required by the boundary-layer model. Separation is
identified when the calculated wall shear stress approaches zero and
changes sign,

\begin{equation}
    \tau_w \rightarrow 0,
\end{equation}

followed by

\begin{equation}
    \tau_w<0.
\end{equation}

% ---------------------------------------------------------
\subsection{Shock detection and present limitations}
% ---------------------------------------------------------

The present MOC formulation assumes a continuous and shock-free
supersonic flowfield. Compression characteristics may converge and
eventually intersect. Such characteristic coalescence indicates that a
smooth isentropic solution no longer exists and that a shock may form.

Until a dedicated shock-fitting method is implemented, characteristic
intersection is treated as a solver-failure condition rather than as a
calculated post-shock solution. Therefore, the present model resolves
smooth expansion and compression waves, but it does not yet calculate
the downstream state across a finite-strength shock.

The main assumptions of the present formulation are:

\begin{enumerate}
    \item steady axisymmetric flow;

    \item inviscid core flow;

    \item no swirl;

    \item calorically perfect gas in the first implementation;

    \item smooth supersonic solution without fitted shocks;

    \item boundary-layer effects treated through a separate integral
    model.
\end{enumerate}

% ---------------------------------------------------------
\subsection{Numerical verification}
% ---------------------------------------------------------

Mass-flow conservation is used as a primary numerical diagnostic. The
relative mass-flow residual is

\begin{equation}
    \boxed{
    \varepsilon_{\dot m}
    =
    \left|
    \frac{\dot m_e-\dot m_t}{\dot m_t}
    \right|
    }.
    \label{eq:moc_mass_residual}
\end{equation}

A solution is accepted for optimization only when the mass-flow
residual, local predictor--corrector residuals and mesh-quality criteria
remain below their prescribed tolerances.

Grid convergence is evaluated by successively refining both the initial
characteristic line and the prescribed wall discretization. The
quantities monitored during refinement include

\begin{equation}
    \dot m_e,\qquad
    C_{F,\mathrm{MOC}},\qquad
    p_e(r),\qquad
    M_e(r),\qquad
    \theta_e(r).
\end{equation}

